We mechanise three theorems in Lean 4 {[}LEA{]}. The full development
lives in \texttt{paper/lean/}; we sketch the statements and proof
structure here.

Let \(\Sigma\) denote the set of well-formed JSON objects under a fixed
schema \(S\) in the supported fragment, and \(\Pi\) the set of policies
in the supported policy fragment. Let \(\text{Tok}\) denote the type of
capability tokens, \(\Pi_\text{tok}: \text{Tok} \to \Pi\) the projection
from a token to its constraint set, and
\(\text{Verify}: \text{Tok} \times \mathcal{K} \to \{0, 1\}\) Ed25519
verification under a public key.

\textbf{Theorem 1 (Attenuation Soundness).} \emph{For all
\(p, c : \text{Tok}\), if \(c\) is in the image of
\(\text{attenuate}(p, \cdot)\), then for every constraint kind \(k\) and
every field \(f\), the narrowing relation \(\sqsubseteq_k\) holds
between \(\Pi_\text{tok}(c)|_{k,f}\) and \(\Pi_\text{tok}(p)|_{k,f}\),
and \(c.\text{expires\_at} \le p.\text{expires\_at}\).}

The proof proceeds by case analysis on \(k\). For each constraint kind,
\texttt{attenuate}'s implementation is shown to compute the appropriate
lattice meet (intersection, union, min, max). Because each operation is
monotone in the lattice, \(\sqsubseteq_k\) is preserved by definition.
The lifetime bound follows from the explicit guard in \texttt{attenuate}
that returns \texttt{none} when
\(c.\text{expires\_at} > p.\text{expires\_at}\). The Lean development is
approximately 105 lines and compiles cleanly under Mathlib v4.10.0.

\textbf{Theorem 2 (Gate Soundness).} \emph{For all \(t : \text{Tok}\),
all tools \(T\), and all argument bindings \(a\), if
\(\text{Gate.check}(t, T, a) = \text{ALLOW}\) under trusted-issuer set
\(\mathcal{K}\), then: (i) \(t.\text{tool} = T\); (ii) every constraint
of the kinds modelled in the development --- \texttt{allowed\_values},
\texttt{forbidden\_values}, \texttt{max\_value}, \texttt{min\_value},
and \texttt{required\_present} --- in \(\Pi_\text{tok}(t)\) is satisfied
by \(a\); (iii) the signature on \(t\) verifies under some
\(k \in \mathcal{K}\); (iv)
\(t.\text{token\_id} = \text{cap:}\Vert\text{SHA256}(\text{canon}(t.\text{body}))\).}

The proof is a direct case analysis on the gate checks. Each check, if
it returned ALLOW, established the corresponding postcondition. The
pre-image analysis is mechanical. The Lean development is approximately
95 lines and compiles cleanly.

\textbf{Scope of the mechanised constraint vocabulary.} The Lean
\texttt{Policy} models the five constraint kinds named in clause (ii).
The reference implementation's runtime gate enforces four further kinds
--- \texttt{starts\_with}, \texttt{forbidden\_field\_combinations},
dot-delimited \texttt{agent\_id} descendant scope, and a revocation
denylist --- together with token expiry and issuer pinning. These are
covered by the implementation's test suite but are \textbf{not} part of
the Lean model; extending the mechanisation to them is future work. We
therefore state Theorem 2 over the modelled vocabulary and do not claim
mechanised soundness for the four runtime-only kinds.

\textbf{Theorem 3 (Policy-Proof Composition).} \emph{Let \(S\) be a
schema, \(P\) a policy, and \(\rho\) a proof artifact for \((S, P)\)
with \(\rho.\text{status} = \text{PROVEN}\). Then for every argument
binding \(a\) that conforms to \(S\) and is accepted by the gate under a
token \(t\) ---
i.e.~\(\text{Gate.check}(t, t.\text{tool}, a) = \text{ALLOW}\) --- both:
(i) \(a\) satisfies \(P\); and (ii) \(a\) satisfies \(t\)'s own
constraints.}

What the mechanised proof establishes, precisely. Clause (i) is
discharged by an explicit \textbf{prover-soundness axiom} standing for
the Z3 oracle: it states that a \texttt{PROVEN} artifact for \((S, P)\)
certifies \(\Sigma \cap \neg P = \emptyset\) --- every schema-conforming
assignment satisfies \(P\) --- so \(a\)'s schema conformance yields
\(a \models P\) directly. This axiom is \textbf{load-bearing} in the
proof of clause (i) (it is applied, not merely recorded). Clause (ii) is
discharged independently from gate soundness (Theorem 2). What is
deliberately \emph{not} mechanised: that the proof \(\rho\) pertains to
this tool's \((S, P)\) and that \(t\) cites it --- the
\(\text{policy\_proof\_hash} \leftrightarrow\) grammar/policy-hash
binding --- is enforced operationally by the gate's strict proof mode
(\texttt{\_check\_proof\_binding} in the reference implementation) and
is taken as a hypothesis of the theorem rather than re-derived from
hashes in Lean. This is the design intent: the SMT solver (via the
axiom), not Lean, certifies policy completeness; the runtime gate checks
the citation. Consequently the mechanised result is strongest when
strict proof mode is enabled --- the default gate mode treats
\texttt{policy\_proof\_hash} as informational and does not enforce the
binding. The Lean development for this theorem is approximately 130
lines. The full development --- the shared data model plus Theorems 1--3
--- is approximately 430 lines of Lean 4 and compiles under Mathlib
v4.10.0 with no \texttt{sorry}s.

We are honest about what the Lean development does and does not
establish. It mechanises the \emph{attenuation and gate logic at the
level of the abstract data model}. It does not mechanise the bit-level
correctness of the Ed25519 implementation (we rely on Erbsen et al.'s
verified field-arithmetic for Curve25519/Ed25519 {[}erbsen2019simple{]}
and the audited \texttt{cryptography} library {[}PYCA{]}) nor the
bit-level correctness of Z3 (we treat the SMT solver as a trusted oracle
whose \texttt{unsat} results we accept; in practice a \texttt{PROVEN}
result from Z3 should be cross-validated against a second
implementation, a step we have added to the reference codebase).

The mechanisation depends on three explicit cryptographic-oracle
assumptions, all standard. (i) \textbf{Ed25519 unforgeability under
chosen-message attack:} the only way to produce a valid signature over a
body is to know the corresponding private key. Theorem 2's conclusion
that the verifying public key is in the trusted set rests on this. (ii)
\textbf{SHA-256 collision resistance:} no two distinct canonical bodies
hash to the same \texttt{token\_id}. Theorem 2's identifier-binding
check rests on this. (iii) \textbf{\texttt{canonBody} determinism:} the
canonical-JSON serialisation produces the same byte sequence for every
logically-equivalent body across implementations and Python versions.
Our canonicalisation uses sorted keys, no whitespace, UTF-8 encoding,
and \texttt{ensure\_ascii=False}; we treat this as an oracle property
and recommend deployments cross-validate the canonical encoding against
a second implementation at boot.
